\section{Sample construction and stress-test method}\label{sec:method}

\subsection{Frozen pool and purposive sample}
The sampling window is 1 August 2026 00:00 UTC through 25 August 2026 23:59 UTC. The candidate pool was obtained from the official arXiv API by querying first submissions in the primary categories \texttt{hep-th}, \texttt{gr-qc}, \texttt{math-ph}, \texttt{quant-ph}, \texttt{cond-mat.stat-mech}, \texttt{hep-ph}, and \texttt{nucl-th}. Cross-listed papers were retained only under their primary category, yielding 2,352 distinct records. The complete pool is frozen in \texttt{data/candidate\_pool.jsonl}; its SHA-256 digest and all query parameters are recorded in \texttt{data/sample\_manifest.json}.

The study uses a stratified purposive stress-test sample: two papers were retained in each of five topic strata, field/particle/nuclear theory, gravity/cosmology, mathematical physics, quantum information, and many-body/statistical physics. The design sought adversarial variation in mathematical form, evidence type, source layout, and semantic pressure. Eligibility favored an original theoretical or computational result about a physical or mathematical system and an official arXiv source archive containing TeX; instrumentation-only, experiment-only, observational-fit-only, review, erratum, and general machine-learning architecture papers without a central physical claim were outside the intended stress-test scope.

An earlier protocol proposed deterministic within-stratum random queues using the seed \path{AgtXIv-ScientificClaim-interface-2026-08-25-v1}. The retained decision log does not reconstruct those queues or every pre-selection eligibility decision and exclusion. The study therefore corrects the design classification to ``stratified purposive'' rather than backfilling an unobserved randomization history. The frozen pool, query parameters, and selected-paper artifacts make the August pool and the analyzed set reproducible; they do not make the ten papers a probability sample or support population-representative inference.

\subsection{Claim-bearing unit screening}
Each paper is scanned section by section rather than only through its abstract. A passage enters the candidate set when it makes a declarative statement that can, in principle, be assessed as true or false relative to an identified system, model, dataset, or body of work. Definitions, assumptions, imported results, informal derivations, theorem statements, numerical findings, limitations, and interpretive conclusions are retained. Pure navigation, acknowledgements, bibliographic metadata, and descriptions that assert no checkable content are excluded.

A sentence is not presumed to equal a claim. A claim may occupy a clause, may require several adjacent clauses to preserve a truth-relevant condition, or may be supported by non-contiguous spans. Conversely, a compound sentence is split when its conclusions can vary independently. Every retained record stores exact source text and line ranges before decontextualization.

\subsection{Stress-test record}
The machine-readable records in \texttt{data/claims-<arXiv-id>.jsonl} follow \texttt{data/claim-record-schema.json}. This is an audit format, not the proposed production interface. It deliberately records information needed to test the interface: source spans, atomicity decisions, conditions, quantifiers, modality, knowledge attribution, suitability for \MathClaimIR{}, and any pressure placed on the proposed boundary. Any automated readiness check over these artifacts verifies declared artifact completeness and internal consistency only; it does not verify random selection, semantic admission, scientific truth, or universality.

\subsection{Two-pass review}
Pass one extracts source-grounded candidates without changing the interface to accommodate them. Pass two reopens the source and checks exact text, omitted conditions, over-splitting, under-splitting, knowledge attribution, and mathematical-normalization suitability. Cross-paper failures are then assigned to one of five destinations: the \ScientificClaim{} core, an optional annotation, a relation, an assessment, or a query/view object. A field enters the core only when it is identity-bearing, cannot be reconstructed without loss, and is required by at least two papers from different strata.
